\chapter{Local Fields}
\begin{center}\itshape
The field $\Q_2$ has exactly seven quadratic extensions, and $\Q_5$ has three.\\
Every finite extension of $\Q_p$ is generated by a root of a polynomial you can write down by hand,\\
and for each degree there are only finitely many of them.
\end{center}

In 1897 Kurt Hensel, a former student of Kronecker in Berlin, had a strange idea. Weierstrass and Puiseux had taught
everyone to study an algebraic function near a point by expanding it in a power series in $(x-a)$. Dedekind and Weber had
just shown that algebraic numbers and algebraic functions obey the same laws. So why not expand a \emph{number} ``near a
prime $p$'', as a power series in $p$? The result was the $p$-adic numbers of Chapter~\ref{ch:padic}, where $-1=\sum_{k\ge0}
(p-1)p^k$ and Newton's method finds square roots digit by digit.

For years almost nobody cared, and Hensel did not help his cause. In 1905 he announced a $p$-adic proof that $e$ is
transcendental: the exponential series $\sum p^n/n!$ converges $p$-adically, and he treated its $p$-adic sum as if it were
the real number $e^p$. It is not, and the ``proof'' collapsed. The rescue came from outside. In 1913 J\'ozsef K\"ursch\'ak
wrote down the axioms of an absolute value; in 1916 Alexander Ostrowski proved that $\Q$ has no others than $|\cdot|_\infty$
and the $|\cdot|_p$ (Theorem~\ref{thm:ostrowski}). And in 1920 a young student named Helmut Hasse found Hensel's book in a
second-hand bookshop in G\"ottingen, was captivated, and moved to Marburg to work with him. Within a year he had proved that
a quadratic form has a rational zero if and only if it has a real zero and a $p$-adic zero for every $p$. That theorem is the
subject of Chapter~9, and it turned the $p$-adic numbers from a curiosity into one of the two legs on which number theory
stands.

This chapter studies the finite extensions of $\Q_p$, the \emph{$p$-adic fields}. They are much tamer than number fields.
The absolute value extends in exactly one way, given by a norm. Each extension has two numbers, the ramification index $e$
and the residue degree $f$, with $ef=n$ exactly, because there is only one prime. The unramified extensions are cyclotomic,
one for each degree. The totally ramified ones are generated by roots of Eisenstein polynomials. And thanks to a beautiful
lemma of Marc Krasner (1946), a polynomial that is close enough to another one generates the same field, so that each
degree has only finitely many extensions. Finally we come back to number fields: completing $K$ at each prime above $p$
splits $K\otimes\Q_p$ into a product of $p$-adic fields, and the global identity $\sum e_if_i=n$ of Chapter~2 reappears as a
sum of local degrees.

\begin{reminder}[from Chapters 1, 2 and 11, and Volume 6]
\emph{$p$-adic numbers} (Chapter~\ref{ch:padic}; Volume~6, Theorem~2.10). $\Q_p$ is the completion of $\Q$ for
$|x|_p=p^{-v_p(x)}$; $|\cdot|_p$ is \emph{ultrametric}: $|x+y|\le\max(|x|,|y|)$, with equality when $|x|\ne|y|$
(Proposition~\ref{prop:ultra}). $\Z_p=\{|x|_p\le1\}$ is compact, $\Z$ is dense in it, and its units are the $x$ with $|x|_p=1$
(Proposition~\ref{prop:Zp}).
\emph{Hensel's lemma} (Theorem~\ref{thm:henselnewton}): if $|f(a)|<|f'(a)|^2$, Newton's method converges to a root; its proof uses only
the ultrametric inequality and completeness, so it holds in every complete ultrametric field. The $p-1$ roots of unity of
$\Z_p$ (Corollary~\ref{cor:teich}). The square classes: $\Q_p^\times/(\Q_p^\times)^2$ has $4$ elements for odd $p$, represented by
$1,u,p,up$ with $u$ a non-square unit, and $8$ for $p=2$, represented by $\pm1,\pm2,\pm5,\pm10$ (Corollary~\ref{cor:squares} and Exercise~\ref{ex:sqclasses}).
\emph{Chapter~2}: unique factorisation of ideals, $v_P$, and $\sum e_if_i=n$.
\end{reminder}

\section{Absolute values on number fields}

An \emph{absolute value} on a field $F$ is a map $|\cdot|:F\to\R_{\ge0}$ with $|x|=0\iff x=0$, $|xy|=|x||y|$ and
$|x+y|\le|x|+|y|$. It is \emph{non-archimedean} (or ultrametric) if $|x+y|\le\max(|x|,|y|)$, and \emph{trivial} if $|x|=1$
for all $x\ne0$. Two absolute values are \emph{equivalent} if $|\cdot|_2=|\cdot|_1^s$ for some $s>0$; equivalent absolute values
define the same topology and the same completion.

Let $K$ be a number field, $\mathcal O=\mathcal O_K$. For a prime $P$ of $\mathcal O$, every $x\in K^\times$ has a valuation
$v_P(x)\in\Z$, the exponent of $P$ in the fractional ideal $x\mathcal O$ (Chapter~2, Corollary~\ref{cor:divides}). Since
$(xy)=(x)(y)$ and $(x+y)\subseteq(x)+(y)$, we have $v_P(xy)=v_P(x)+v_P(y)$ and $v_P(x+y)\ge\min(v_P(x),v_P(y))$. Hence
\[
|x|_P=N(P)^{-v_P(x)}
\]
is a non-archimedean absolute value on $K$. For each embedding $\sigma:K\to\C$, $x\mapsto|\sigma(x)|$ is an archimedean one,
and conjugate embeddings give the same.

\begin{theorem}[Ostrowski's theorem for number fields, non-archimedean part]\label{thm:ostrowskiK}
Every non-trivial non-archimedean absolute value on $K$ is equivalent to $|\cdot|_P$ for exactly one prime $P$ of
$\mathcal O_K$.
\end{theorem}
\begin{proof}
Let $|\cdot|$ be one. Then $|n|=|1+\dots+1|\le1$ for $n\in\Z$. If $\alpha\in\mathcal O$ satisfies $\alpha^d=-\sum_{i<d}a_i
\alpha^i$ with $a_i\in\Z$ and we had $|\alpha|>1$, then $|\alpha|^d\le\max_i|\alpha|^i<|\alpha|^d$; so $|\alpha|\le1$ on $\mathcal
O$. The set $P=\{\alpha\in\mathcal O:|\alpha|<1\}$ is an ideal (ultrametric inequality, and $|\beta\alpha|\le|\alpha|$ for
$\beta\in\mathcal O$), and it is prime ($|\alpha\beta|<1$ forces $|\alpha|<1$ or $|\beta|<1$). It is not zero: otherwise
$|\alpha|=1$ for all $\alpha\in\mathcal O\setminus\{0\}$, and $|\cdot|$ would be trivial on $K=\operatorname{Frac}\mathcal O$.

Let $y\in K^\times$ with $v_P(y)=0$. Write $y\mathcal O=AB^{-1}$ with $A,B$ ideals not divisible by $P$. Pick $b\in B\setminus P$
(possible since $B\not\subseteq P$, Chapter~2, Corollary~\ref{cor:divides}(b)); then $by\in bAB^{-1}\subseteq A\subseteq\mathcal O$,
and $by\notin P$ since $v_P(by)=0$. So $|by|=|b|=1$ and $|y|=1$. Now choose $\pi\in P\setminus P^2$, so $v_P(\pi)=1$, and for
$x\in K^\times$ apply this to $y=x\pi^{-v_P(x)}$: $|x|=|\pi|^{v_P(x)}$. With $s=\log|\pi|/\log N(P)^{-1}>0$ this is $|x|_P^{\,s}$.

Uniqueness: if $P\ne P'$, pick $x\in P\setminus P'$ (maximal ideals are not contained in one another); then $|x|_P<1$ and
$|x|_{P'}=1$, so $|\cdot|_P$ and $|\cdot|_{P'}$ are not equivalent.
\end{proof}

\begin{remark}
The archimedean absolute values of $K$ are, up to equivalence, the $|\sigma(\cdot)|$, one for each real embedding and one
for each pair of complex embeddings. The proof rests on Ostrowski's other theorem, that $\R$ and $\C$ are the only fields
complete for an archimedean absolute value (Neukirch, Chapter~II, Theorem~4.2, or Cassels, \emph{Local Fields}, Chapter~3).
The equivalence classes of absolute values of $K$ are called its \emph{places}: finite places are primes, infinite places
are embeddings up to conjugation, and there are $r_1+r_2$ of them.
\end{remark}

We normalise: $|x|_P=N(P)^{-v_P(x)}$ as above, $|x|_\sigma=|\sigma(x)|$ for a real embedding, and $|x|_\sigma=|\sigma(x)|^2$
for a complex one (not an absolute value, because of the square, but the right normalisation).

\begin{proposition}[the product formula]\label{prop:productformula}
For every $x\in K^\times$, $\displaystyle\prod_v|x|_v=1$, the product over all places of $K$.
\end{proposition}
\begin{proof}
Only finitely many $v_P(x)$ are non-zero. By multiplicativity of the norm of ideals (Chapter~2, Proposition~\ref{prop:normideal}),
extended to fractional ideals, $\prod_PN(P)^{v_P(x)}=N(x\mathcal O)=|N_{K/\Q}(x)|$, so $\prod_P|x|_P=|N(x)|^{-1}$. And
$\prod_{\text{infinite}}|x|_v=\prod_{\sigma\text{ real}}|\sigma(x)|\prod_{\text{pairs}}|\sigma(x)|^2=|N(x)|$.
\end{proof}

For $K=\Q$ it says $|x|_\infty\prod_p|x|_p=1$, which is unique factorisation in disguise: $12=2^2\cdot3$ and indeed
$12\cdot\frac14\cdot\frac13=1$.

\section{Vector spaces over a complete field, and Hensel's factorisation}

Let $F$ be a field complete for a non-trivial absolute value $|\cdot|$. A \emph{norm} on an $F$-vector space $V$ is a map
$\|\cdot\|:V\to\R_{\ge0}$ with $\|v\|=0\iff v=0$, $\|av\|=|a|\|v\|$ and $\|v+w\|\le\|v\|+\|w\|$. If $e_1,\dots,e_n$ is a basis,
$|\sum a_ie_i|_\infty=\max_i|a_i|$ is a norm, the \emph{sup norm}, and $V$ is complete for it, coordinate by coordinate.

\begin{theorem}\label{thm:normequiv}
On a finite-dimensional vector space over a complete field $F$, all norms are equivalent: for two norms there are $c,C>0$
with $c\|v\|_1\le\|v\|_2\le C\|v\|_1$. In particular $V$ is complete for every norm.
\end{theorem}
\begin{proof}
It suffices to compare a norm $\|\cdot\|$ with the sup norm in a basis. One side is easy: $\|\sum a_ie_i\|\le\sum|a_i|\|e_i\|
\le C|v|_\infty$ with $C=\sum\|e_i\|$. For the other side we use induction on $n=\dim V$; for $n=1$, $\|ae_1\|=|a|\|e_1\|$.
Suppose there is no $c>0$ with $\|v\|\ge c|v|_\infty$. Then there are $v_k$ with $\|v_k\|<\frac1k|v_k|_\infty$. Divide $v_k$ by
a coordinate of largest absolute value; the new $v_k$ has that coordinate equal to $1$, all coordinates of absolute value at
most $1$, and $\|v_k\|<\frac1k$. Infinitely many $k$ share the same position $j$ of that coordinate; keeping only those,
$v_k=e_j+w_k$ with $w_k\in W=\operatorname{span}(e_i:i\ne j)$. Now $\|w_k-w_l\|=\|v_k-v_l\|<\frac1k+\frac1l$, so $(w_k)$ is Cauchy for
$\|\cdot\|$ on $W$. By induction, $\|\cdot\|$ and the sup norm are equivalent on $W$, which is complete for the sup norm; so
$w_k\to w\in W$ for the sup norm, hence for $\|\cdot\|$. Then $\|e_j+w\|=\lim\|v_k\|=0$, so $e_j=-w\in W$, which is absurd.
Completeness follows: a Cauchy sequence for $\|\cdot\|$ is Cauchy for the sup norm, hence converges.
\end{proof}

From now on $F$ is a field complete for a non-archimedean absolute value whose value group $|F^\times|$ is discrete, hence
cyclic: $|F^\times|=\langle|\pi_F|\rangle$. We call such an $F$ a \emph{complete discretely valued field}. Its
\emph{valuation ring} is $\mathcal O_F=\{|x|\le1\}$, with the unique maximal ideal $\pi_F\mathcal O_F=\{|x|<1\}$ (an element of
absolute value $1$ is invertible in $\mathcal O_F$) and \emph{residue field} $k_F=\mathcal O_F/\pi_F\mathcal O_F$. The example to
keep in mind is $F=\Q_p$, $\mathcal O_F=\Z_p$, $\pi_F=p$, $k_F=\mathbb F_p$. Note that $\mathcal O_F$ is integrally closed: if
$|x|>1$ and $x^d=-\sum_{i<d}a_ix^i$ with $a_i\in\mathcal O_F$, the right side has absolute value $<|x|^d$. Reduction modulo
$\pi_F$ is written $a\mapsto\bar a$, on elements and on polynomials.

\begin{theorem}[Hensel's lemma, factorisation form]\label{thm:henselfact}
Let $f\in\mathcal O_F[X]$ have degree $d$, and let $g_0,h_0\in\mathcal O_F[X]$, with $g_0$ monic of degree $m$ and $\deg h_0\le
d-m$, such that $\bar f=\bar g_0\bar h_0$ and $\bar g_0,\bar h_0$ are coprime in $k_F[X]$. Then $f=gh$ with $g,h\in\mathcal
O_F[X]$, $g$ monic of degree $m$, $\bar g=\bar g_0$ and $\bar h=\bar h_0$.
\end{theorem}
\begin{proof}
Write $\pi=\pi_F$ and choose $a,b\in\mathcal O_F[X]$ with $\bar a\bar g_0+\bar b\bar h_0=1$ (Bezout in $k_F[X]$). We build
$g_k$ monic of degree $m$ and $h_k$ of degree $\le d-m$, with $g_k\equiv g_0$, $h_k\equiv h_0$ and $f\equiv g_kh_k\pmod{\pi^{k+1}}$,
starting from $k=0$. Given them, write $f-g_kh_k=\pi^{k+1}c$ with $c\in\mathcal O_F[X]$ of degree $\le d$. In $k_F[X]$, divide
$\bar b\bar c$ by the monic $\bar g_0$: $\bar b\bar c=\bar q\bar g_0+\bar v$ with $\deg\bar v<m$. Then
\[
\bar c=\bar c(\bar a\bar g_0+\bar b\bar h_0)=\bar u\bar g_0+\bar v\bar h_0,\qquad\bar u=\bar a\bar c+\bar q\bar h_0.
\]
Since $\bar u\bar g_0=\bar c-\bar v\bar h_0$ has degree $\le d$, $\deg\bar u\le d-m$. Lift $\bar u,\bar v$ to $u,v\in\mathcal
O_F[X]$ of the same degrees and put $g_{k+1}=g_k+\pi^{k+1}v$, $h_{k+1}=h_k+\pi^{k+1}u$. Then $g_{k+1}$ is monic of degree $m$
($\deg v<m$), $\deg h_{k+1}\le d-m$, and
\[
f-g_{k+1}h_{k+1}=\pi^{k+1}(c-vh_k-ug_k)-\pi^{2k+2}uv\equiv\pi^{k+1}(c-vh_0-ug_0)\equiv0\pmod{\pi^{k+2}}.
\]
The coefficients of $g_k$ and $h_k$ form Cauchy sequences in the complete ring $\mathcal O_F$; their limits $g,h$ satisfy $f=gh$.
\end{proof}

\begin{corollary}\label{cor:integralcoef}
Let $f\in F[X]$ be monic and irreducible, with $f(0)\in\mathcal O_F$. Then $f\in\mathcal O_F[X]$.
\end{corollary}
\begin{proof}
Suppose not, and let $f=X^n+\dots+a_0$. Multiply $f$ by a power of $\pi$ to get $\tilde f=\pi^sf\in\mathcal O_F[X]$ with some
coefficient of absolute value $1$; here $s\ge1$. Let $r$ be the largest index with $|\tilde a_r|=1$. The leading coefficient
$\pi^s$ and the constant term $\pi^sa_0$ are in $\pi\mathcal O_F$, so $0<r<n$. Then $\bar{\tilde f}=\bar{\tilde a}_r\cdot\bar g_0$
with $\bar g_0$ monic of degree $r$, and $\bar{\tilde a}_r$, a non-zero constant, is coprime to it. Theorem~\ref{thm:henselfact}
gives $\tilde f=gh$ with $\deg g=r$, hence $\deg h=n-r$, both positive: $f$ is reducible. Contradiction.
\end{proof}

\section{Extending the absolute value}

\begin{theorem}\label{thm:extabs}
Let $F$ be a complete discretely valued field and $L/F$ a finite extension of degree $n$. Then
\[
|x|_L=|N_{L/F}(x)|^{1/n}
\]
is an absolute value on $L$ extending $|\cdot|$, it is the only one, and $L$ is complete for it. Its valuation ring $\mathcal
O_L=\{|x|_L\le1\}$ is the integral closure of $\mathcal O_F$ in $L$.
\end{theorem}
\begin{proof}
\emph{Integrality.} Let $x\in L$ with minimal polynomial $m_x$ of degree $d$ over $F$. Then $N_{L/F}(x)=\bigl((-1)^dm_x(0)\bigr)^{n/d}$
(the characteristic polynomial of multiplication by $x$ is $m_x^{n/d}$, by the proof of Chapter~1, Proposition~\ref{prop:charpoly}, which works over any field). So $|x|_L\le1$
iff $m_x(0)\in\mathcal O_F$ iff $m_x\in\mathcal O_F[X]$ (Corollary~\ref{cor:integralcoef}). And $m_x\in\mathcal O_F[X]$ iff $x$
is integral over $\mathcal O_F$: if $x$ is a root of a monic $g\in\mathcal O_F[X]$, the roots of $m_x$ (in a splitting field)
are roots of $g$, so the coefficients of $m_x$ are integral over $\mathcal O_F$ and lie in $F$, hence in the integrally closed
$\mathcal O_F$.

\emph{It is an absolute value.} $|x|_L=0$ iff $x=0$, and $|xy|_L=|x|_L|y|_L$ since the norm is multiplicative. For the
ultrametric inequality let $|x|_L\le|y|_L$, $y\ne0$. Then $z=x/y$ has $|z|_L\le1$, so $m_z\in\mathcal O_F[X]$, and $1+z$ has
minimal polynomial $m_z(X-1)\in\mathcal O_F[X]$; so $|1+z|_L\le1$ and $|x+y|_L\le|y|_L$. On $F$, $N_{L/F}(a)=a^n$, so $|a|_L=|a|$.

\emph{Uniqueness and completeness.} Any absolute value extending $|\cdot|$ is a norm on the $F$-vector space $L$. By
Theorem~\ref{thm:normequiv}, two of them satisfy $c|x|_1\le|x|_2\le C|x|_1$; applying this to $x^k$ and taking $k$-th roots,
$c^{1/k}|x|_1\le|x|_2\le C^{1/k}|x|_1$, and $k\to\infty$ gives $|x|_1=|x|_2$. Completeness is part of
Theorem~\ref{thm:normequiv}.
\end{proof}

\begin{corollary}\label{cor:galoisinv}
Fix an algebraic closure $\overline{\Q}_p$. There is exactly one absolute value on $\overline{\Q}_p$ extending $|\cdot|_p$,
and it satisfies $|\sigma(x)|=|x|$ for every automorphism $\sigma$ of $\overline{\Q}_p$ and every $x$; in particular conjugate
elements over a $p$-adic field have the same absolute value.
\end{corollary}
\begin{proof}
Every $x$ lies in the finite extension $\Q_p(x)$, and the absolute values defined there are compatible by uniqueness. For
$\sigma$, $x\mapsto|\sigma(x)|$ is again an absolute value extending $|\cdot|_p$, so it is $|\cdot|$.
\end{proof}

\begin{think}
Is $\Q_p$ algebraically closed? (No: $X^2-p$ has no root, since $|\sqrt p|=p^{-1/2}$ is not a power of $p$.) Is
$\overline{\Q}_p$, which has infinite degree, still complete? Guess before reading on. (It is not. Its completion $\C_p$ is
complete and, surprisingly, algebraically closed again; see Gouv\^ea, \S5.7. In this book we never need $\C_p$.)
\end{think}

\section{Ramification index and residue degree}

Let $F$ be complete discretely valued and $L/F$ of degree $n$. Since $|L^\times|=|N_{L/F}(L^\times)|^{1/n}\subseteq
|F^\times|^{1/n}$, the value group of $L$ is discrete, so $L$ is complete discretely valued too. Write
$|L^\times|=\langle|\pi_L|\rangle$, with $\pi_L$ a \emph{uniformiser} of $L$.

\begin{definition}
The \emph{ramification index} of $L/F$ is $e=[\,|L^\times|:|F^\times|\,]$, so that $|\pi_F|=|\pi_L|^e$ and $\pi_F\mathcal O_L=
\pi_L^e\mathcal O_L$. The \emph{residue degree} is $f=[k_L:k_F]$. The extension is \emph{unramified} if $e=1$ and
\emph{totally ramified} if $e=n$.
\end{definition}

\begin{theorem}\label{thm:efn}
$ef=n$. More precisely, if $\omega_1,\dots,\omega_f\in\mathcal O_L$ reduce to a basis of $k_L$ over $k_F$, then the $ef$
elements $\omega_i\pi_L^j$ ($1\le i\le f$, $0\le j<e$) form a basis of $\mathcal O_L$ over $\mathcal O_F$ and of $L$ over $F$.
\end{theorem}
\begin{proof}
\emph{Step 1.} If $\omega_1,\dots,\omega_r\in\mathcal O_L$ have residues linearly independent over $k_F$, then $|\sum a_i\omega_i|=
\max|a_i|$ for $a_i\in F$. Indeed, dividing by an $a_{i_0}$ of largest absolute value we may assume $\max|a_i|=1$; then
$\sum\bar a_i\bar\omega_i\ne0$, so $\sum a_i\omega_i$ has absolute value $1$. In particular such $\omega_i$ are linearly
independent over $F$, so $r\le n$ and $f$ is finite.

\emph{Step 2: independence.} Let $\sum_{i,j}a_{ij}\omega_i\pi_L^j=0$, and put $s_j=\sum_ia_{ij}\omega_i$, so that
$|s_j|=\max_i|a_{ij}|\in|F^\times|\cup\{0\}$ by Step~1. The non-zero terms $s_j\pi_L^j$ have absolute values in the distinct
cosets $|\pi_L|^j|F^\times|$, $0\le j<e$, so they are pairwise different, and by the ultrametric inequality (with equality for
different absolute values) a sum of such terms is non-zero unless all terms vanish. Hence all $a_{ij}=0$, and $ef\le n$.

\emph{Step 3: spanning.} Let $M=\sum_{i,j}\mathcal O_F\omega_i\pi_L^j$. Every $x\in\mathcal O_L$ has $\bar x=\sum\bar a_i
\bar\omega_i$, so $\mathcal O_L=\sum_i\mathcal O_F\omega_i+\pi_L\mathcal O_L$; iterating $e$ times, $\mathcal O_L=M+\pi_L^e\mathcal
O_L=M+\pi_F\mathcal O_L$. So $x=m_0+\pi_Fx_1$, $x_1=m_1+\pi_Fx_2$, and so on, with $m_k\in M$: $x=\sum_{k<K}\pi_F^km_k+\pi_F^Kx_K$.
The coefficients of $\sum_k\pi_F^km_k$ in the $\omega_i\pi_L^j$ converge in $\mathcal O_F$, and $\pi_F^Kx_K\to0$, so $x\in M$.
Thus $\mathcal O_L=M$, free with basis $(\omega_i\pi_L^j)$. Every $x\in L$ has $\pi_F^kx\in\mathcal O_L$ for large $k$, so this is
also a basis of $L$, and $n=ef$.
\end{proof}

Theorems~\ref{thm:extabs} and~\ref{thm:efn} apply to every finite extension of a finite extension of $\Q_p$, since each such
field is complete discretely valued, with a finite residue field. We call these fields \emph{$p$-adic fields}. For a tower
$F\subseteq L\subseteq M$, both $e$ and $f$ are multiplicative (indices of subgroups and degrees of fields multiply).

\section{Unramified and totally ramified extensions}

In this section $F$ is a $p$-adic field with residue field $\mathbb F_q$, and all extensions lie in $\overline{\Q}_p$.

\begin{theorem}[unramified extensions]\label{thm:unram}
For each $f\ge1$ there is exactly one unramified extension of $F$ of degree $f$, namely $F_f=F(\zeta_{q^f-1})$. It is Galois
over $F$, and reduction gives an isomorphism $\operatorname{Gal}(F_f/F)\cong\operatorname{Gal}(\mathbb F_{q^f}/\mathbb F_q)$,
cyclic of order $f$, generated by the \emph{Frobenius} $\zeta\mapsto\zeta^q$ on roots of unity of order prime to $p$.
\end{theorem}
\begin{proof}
\emph{Roots of unity.} Let $L$ be a $p$-adic field with residue field $\mathbb F_{q^f}$, and $N=q^f-1$. The polynomial $X^N-1$
has $N$ distinct roots in $\mathbb F_{q^f}$, all simple ($N$ is prime to $p$). By Hensel's lemma (Newton's form, valid in $\mathcal
O_L$), each lifts to a unique root in $\mathcal O_L$; so $L$ contains $N$ distinct $N$-th roots of unity, with distinct
residues, forming a cyclic group (Volume~2, Theorem~9.4) generated by some $\zeta_N$.

\emph{Existence.} Let $\bar g\in\mathbb F_q[X]$ be irreducible of degree $f$ and $g\in\mathcal O_F[X]$ a monic lift; $g$ is
irreducible over $F$, since a factorisation into monic factors would have factors in $\mathcal O_F[X]$ (Corollary~\ref{cor:integralcoef})
and would reduce to a factorisation of $\bar g$. In $L=F[X]/(g)$, of degree $f$, the class of $X$ is integral and reduces to a
root of $\bar g$; so $k_L\supseteq\mathbb F_{q^f}$, $f(L/F)\ge f=[L:F]$, and Theorem~\ref{thm:efn} forces $e=1$: $L$ is
unramified of degree $f$.

\emph{Uniqueness.} If $L$ is unramified of degree $f$, its residue field is $\mathbb F_{q^f}$, so $\zeta_N\in L$, and the residue
field of $F(\zeta_N)\subseteq L$ already contains the residues of all $N$-th roots of unity, that is all of $\mathbb F_{q^f}$. So
$[F(\zeta_N):F]\ge f$, and $L=F(\zeta_N)$, the field generated over $F$ by all $N$-th roots of unity in $\overline{\Q}_p$.

\emph{Galois group.} $F_f$ contains all roots of $X^N-1$, so it is the splitting field of this polynomial, and Galois.
An automorphism $\sigma$ preserves $|\cdot|$ (Corollary~\ref{cor:galoisinv}), so it maps $\mathcal O_{F_f}$ and its maximal ideal to
themselves and induces $\bar\sigma\in\operatorname{Gal}(\mathbb F_{q^f}/\mathbb F_q)$. If $\bar\sigma=1$, then $\sigma(\zeta_N)$ is an
$N$-th root of unity with the same residue as $\zeta_N$, hence equal to it, and $\sigma=1$. Both groups have order $f$, so
$\sigma\mapsto\bar\sigma$ is an isomorphism; the Frobenius $\bar x\mapsto\bar x^q$ lifts to $\sigma$ with $\sigma(\zeta)\equiv
\zeta^q$, hence $\sigma(\zeta)=\zeta^q$, for every $N$-th root of unity $\zeta$.
\end{proof}

\begin{theorem}[Eisenstein]\label{thm:eisensteinlocal}
\begin{enumerate}[label=(\alph*)]
\item If $g=X^e+a_{e-1}X^{e-1}+\dots+a_0\in\mathcal O_F[X]$ is Eisenstein ($|a_i|<1$ for all $i$, $|a_0|=|\pi_F|$), it is
irreducible, and $F(\pi)$ is totally ramified of degree $e$ for every root $\pi$, which is a uniformiser.
\item Conversely, if $L/F$ is totally ramified of degree $e$, every uniformiser $\pi_L$ has an Eisenstein minimal polynomial,
and $L=F(\pi_L)$.
\end{enumerate}
\end{theorem}
\begin{proof}
(a) Let $\pi$ be a root and $d=[F(\pi):F]\le e$. If $|\pi|\ge1$ then $|\pi^e|>|a_i\pi^i|$ for $i<e$, and $g(\pi)\ne0$; so
$|\pi|<1$, then $|a_i\pi^i|<|\pi_F|$ for $1\le i<e$ and $|a_0|=|\pi_F|$, so $g(\pi)=0$ forces $|\pi^e|=|a_0|=|\pi_F|$. Thus
$|\pi|=|\pi_F|^{1/e}$, the ramification index of $F(\pi)/F$ is at least $e$, and $e\le e(F(\pi)/F)\le d\le e$. So $d=e$, $g$ is
the minimal polynomial of $\pi$ (irreducible), $F(\pi)/F$ is totally ramified, and $|\pi|$ generates $|F(\pi)^\times|$.

(b) $\pi_L$ has degree $\ge e(F(\pi_L)/F)=e$, so $L=F(\pi_L)$; its minimal polynomial $g$ has coefficients in $\mathcal O_F$
(Theorem~\ref{thm:extabs}). In a Galois closure, all conjugates of $\pi_L$ have absolute value $|\pi_L|<1$, so every
coefficient of $g$, a sum of products of conjugates, has absolute value $<1$; and $|a_0|=|N_{L/F}(\pi_L)|=|\pi_L|^e=|\pi_F|$.
\end{proof}

\begin{proposition}\label{prop:tower}
Let $L/F$ be an extension of $p$-adic fields with invariants $e,f$. Then $T=F(\zeta_{q^f-1})\subseteq L$, $T/F$ is unramified of
degree $f$, $L/T$ is totally ramified of degree $e$, and $L=T(\pi_L)$ with an Eisenstein polynomial over $\mathcal O_T$.
\end{proposition}
\begin{proof}
$L$ has residue field $\mathbb F_{q^f}$, so it contains $\zeta_{q^f-1}$ (first step of the proof of Theorem~\ref{thm:unram}), and
$T\subseteq L$ is unramified of degree $f$ by that theorem. By multiplicativity, $f(L/T)=1$ and $[L:T]=e$, so $L/T$ is totally
ramified, and Theorem~\ref{thm:eisensteinlocal}(b) applies.
\end{proof}

As a first harvest, here are all quadratic extensions. Since $p$-adic fields have characteristic $0$, every quadratic extension of
$\Q_p$ is $\Q_p(\sqrt d)$, and $\Q_p(\sqrt d)=\Q_p(\sqrt{d'})$ iff $d/d'$ is a square. So the quadratic extensions correspond to the
non-trivial square classes: \emph{three for odd $p$, and seven for $p=2$} (Figure~\ref{fig:quadratic2} and Exercise~\ref{ex:seven}).
For odd $p$, exactly one of the three is unramified, $\Q_p(\sqrt u)=\Q_p(\zeta_{p^2-1})$, and the other two are $\Q_p(\sqrt p)$ and
$\Q_p(\sqrt{up})$, both Eisenstein.

\begin{figure}[ht]
\centering
\fig{ch06-quadratic2}
\caption{The seven quadratic extensions of $\Q_2$. Only $\Q_2(\sqrt5)$ is unramified; the others are generated by roots of
Eisenstein polynomials, and the power of $2$ in their discriminant is $2^2$ or $2^3$ (Exercise~\ref{ex:seven}).}
\label{fig:quadratic2}
\end{figure}

\section{Krasner's lemma and finiteness}

\begin{lemma}[Krasner, 1946]\label{lem:krasner}
Let $F$ be a $p$-adic field and $\alpha,\beta\in\overline{\Q}_p$. Let $\alpha=\alpha_1,\alpha_2,\dots,\alpha_d$ be the conjugates of
$\alpha$ over $F$. If $|\beta-\alpha|<|\alpha-\alpha_i|$ for all $i\ge2$, then $F(\alpha)\subseteq F(\beta)$.
\end{lemma}
\begin{proof}
If not, then $\alpha$ is not fixed by the automorphisms of $\overline{\Q}_p$ over $F(\beta)$ (Galois theory, Volume~2, Theorem~10.6,
in a finite Galois extension of $F(\beta)$ containing $\alpha$), so some such $\sigma$ has $\sigma(\alpha)=\alpha_i$ with $i\ge2$.
By Corollary~\ref{cor:galoisinv}, $|\beta-\alpha_i|=|\sigma(\beta-\alpha)|=|\beta-\alpha|$. Then
$|\alpha-\alpha_i|\le\max(|\alpha-\beta|,|\beta-\alpha_i|)=|\beta-\alpha|<|\alpha-\alpha_i|$, which is absurd.
\end{proof}

\begin{figure}[ht]
\centering
\fig{ch06-krasner}
\caption{Krasner's lemma: if $\beta$ is closer to $\alpha$ than every other conjugate of $\alpha$, then $\alpha\in F(\beta)$. In
an ultrametric world every point of a disc is its centre, so $\beta$ ``sees'' its conjugates exactly as $\alpha$ does.}
\label{fig:krasner}
\end{figure}

\begin{proposition}[roots depend continuously on the polynomial]\label{prop:continuity}
Let $g\in\mathcal O_F[X]$ be monic, irreducible of degree $e$, with roots $\alpha_1,\dots,\alpha_e$, and $\delta=\min_{i\ne j}
|\alpha_i-\alpha_j|$. If $h\in\mathcal O_F[X]$ is monic of degree $e$ and every coefficient of $h-g$ has absolute value $<\delta^e$,
then $h$ is irreducible and each root $\beta$ of $h$ satisfies $F(\beta)=F(\alpha_i)$ for some $i$.
\end{proposition}
\begin{proof}
$\beta$ is integral, so $|\beta|\le1$ and $|g(\beta)|=|g(\beta)-h(\beta)|<\delta^e$. Since $g(\beta)=\prod_i(\beta-\alpha_i)$, some
$|\beta-\alpha_i|<\delta$, and Krasner gives $F(\alpha_i)\subseteq F(\beta)$. As $[F(\beta):F]\le e=[F(\alpha_i):F]$, the fields
are equal, and $\beta$ has degree $e$: $h$ is its minimal polynomial.
\end{proof}

\begin{theorem}\label{thm:finitelymany}
A $p$-adic field $F$ has only finitely many extensions of each degree $n$ in $\overline{\Q}_p$.
\end{theorem}
\begin{proof}
An extension $L$ of degree $n$ has $ef=n$, and by Proposition~\ref{prop:tower} $L=T(\pi_L)$ where $T=F_f$ depends only on $f$
(Theorem~\ref{thm:unram}) and $\pi_L$ is a root of an Eisenstein polynomial of degree $e$ over $\mathcal O_T$. Fix $f$ and $T$.
The set $\mathcal E$ of Eisenstein polynomials of degree $e$ over $\mathcal O_T$ is compact: it is identified with
$(\pi_T\mathcal O_T)^{e-1}\times\pi_T\mathcal O_T^\times$, and $\mathcal O_T$ is compact, being homeomorphic to $\Z_p^{[T:\Q_p]}$ by
Theorem~\ref{thm:efn}. For $g\in\mathcal E$, Proposition~\ref{prop:continuity} gives an open neighbourhood $U_g$ such that the
roots of every $h\in U_g\cap\mathcal E$ generate one of the at most $e$ fields $T(\alpha)$, $g(\alpha)=0$. Finitely many $U_g$
cover $\mathcal E$, so finitely many fields arise.
\end{proof}

\begin{remark}
Krasner went much further and counted the totally ramified extensions of each degree. Serre found the most elegant form of
the count in 1978, his \emph{mass formula}: if $F$ has residue field $\mathbb F_q$, then
\[
\sum_L q^{-(v_F(d_{L/F})-n+1)}=n,
\]
the sum over the totally ramified extensions $L$ of $F$ of degree $n$ inside $\overline{\Q}_p$, where $v_F(d_{L/F})$ is the
valuation of the discriminant. For $F=\Q_2$ and $n=2$: two fields with $v(d)=2$ and four with $v(d)=3$
(Figure~\ref{fig:quadratic2}), and $2\cdot\frac12+4\cdot\frac14=2$. A direct computation verifies more cases.
\end{remark}

\section{Completions of a number field}

Now we come back to a number field $K$ of degree $n$ and a prime number $p$. Write $K=\Q(\theta)$ with $\theta\in\mathcal O_K$
and minimal polynomial $f$. Over $\Q_p$, $f$ factors into distinct monic irreducibles, $f=f_1\cdots f_g$ (distinct because $f$ is
separable); fix a root $\theta_i\in\overline{\Q}_p$ of $f_i$ and let $L_i=\Q_p(\theta_i)$, $\iota_i:K\to L_i$, $\theta\mapsto\theta_i$.

\begin{theorem}\label{thm:completions}
\begin{enumerate}[label=(\alph*)]
\item The absolute values on $K$ extending $|\cdot|_p$ are the $|x|_i=|\iota_i(x)|$, $1\le i\le g$, and they are pairwise
distinct. $L_i$ is the completion of $K$ for $|\cdot|_i$.
\item There is a bijection $i\leftrightarrow P_i$ with the primes of $\mathcal O_K$ above $p$, such that $|x|_i=
|x|_{P_i}^{1/(e_if_i)}$, where $e_i,f_i$ are the invariants of $P_i$ (Chapter~2).
\item $e(L_i/\Q_p)=e_i$ and $f(L_i/\Q_p)=f_i$. Hence $\deg f_i=e_if_i$, and $\sum_ie_if_i=n$.
\end{enumerate}
In other words, $K\otimes_\Q\Q_p\cong\Q_p[X]/(f)\cong\prod_i L_i=\prod_{P\mid p}K_P$: the $p$-adic shadow of $K$ is the product of
its completions.
\end{theorem}
\begin{proof}
(a) Each $|\cdot|_i$ is an absolute value on $K$ extending $|\cdot|_p$, since $\iota_i$ is an injective ring map fixing $\Q$.
$\iota_i(K)=\Q(\theta_i)$ is dense in $L_i=\Q_p+\Q_p\theta_i+\dots+\Q_p\theta_i^{\deg f_i-1}$, because $\Q$ is dense in $\Q_p$; and $L_i$
is complete (Theorem~\ref{thm:extabs}); so $L_i$ is the completion. Conversely let $|\cdot|$ extend $|\cdot|_p$, and let $\hat K$ be
the completion. It contains the closure $\Q_p$ of $\Q$ and $\theta$, hence the field $\Q_p(\theta)$, which is complete, contains
$K$, and lies in the closure of $K$; so $\hat K=\Q_p(\theta)$. There, $\theta$ is a root of $f$, hence of exactly one $f_i$, and
$\Q_p(\theta)\cong L_i$ by $\theta\mapsto\theta_i$; by uniqueness of the absolute value on $L_i$, $|\cdot|=|\cdot|_i$. If
$|\cdot|_i=|\cdot|_j$, the identity of $K$ extends by continuity to an isomorphism $L_i\to L_j$ sending $\theta_i$ to $\theta_j$, so
these have the same minimal polynomial over $\Q_p$, and $i=j$.

(b) By Theorem~\ref{thm:ostrowskiK}, $|\cdot|_i=|\cdot|_P^{\,s}$ for a unique prime $P$, and $P\mid p$ because $|p|_i<1$.
Conversely, for $P\mid p$ with invariants $e,f$, the absolute value $|\cdot|_P^{1/(ef)}$ extends $|\cdot|_p$, since $|p|_P=N(P)^{-e}=
p^{-ef}$; by (a) it is some $|\cdot|_i$. This gives the bijection, and the formula.

(c) By (b), $|x|_i=p^{-v_P(x)/e_i}$, so $|K^\times|_i=p^{\Z/e_i}$. The value group of $L_i$ is the same: for $y\in L_i^\times$ pick
$x\in K$ with $|y-x|<|y|$; then $|x|=|y|$. So $e(L_i/\Q_p)=e_i$. For the residue field, map $\mathcal O_K\to\mathcal O_{L_i}/\mathfrak
m_{L_i}$. The kernel is $\{v_P(x)\ge1\}=P$. It is onto: given $y\in\mathcal O_{L_i}$, pick $x\in K$ with $|y-x|<1$; then
$v_P(x)\ge0$, so (as in the proof of Theorem~\ref{thm:ostrowskiK}) $x=a/b$ with $a,b\in\mathcal O_K$, $b\notin P$; choose $b'\in\mathcal
O_K$ with $bb'\equiv1\pmod P$ ($\mathcal O_K/P$ is a field) and put $x'=ab'\in\mathcal O_K$. Then $x-x'=\frac{a(1-bb')}{b}$ has
$v_P\ge1$, so $x'\equiv x\equiv y$ modulo $\mathfrak m_{L_i}$. Hence $k_{L_i}\cong\mathcal O_K/P$ and $f(L_i/\Q_p)=f_i$.
Theorem~\ref{thm:efn} gives $\deg f_i=[L_i:\Q_p]=e_if_i$, and summing, $\sum_ie_if_i=\sum_i\deg f_i=n$. Finally the Chinese
remainder theorem gives $\Q_p[X]/(f)\cong\prod_i\Q_p[X]/(f_i)\cong\prod_iL_i$.
\end{proof}

\begin{think}
Compare with Kummer--Dedekind (Chapter~2, Theorem~\ref{thm:kd}), which factors $f$ modulo $p$. Here we factor $f$ over $\Q_p$,
which is finer: the factorisation modulo $p$ of each $f_i$ is $\bar g_i^{\,e_i}$ when $p\nmid[\mathcal O_K:\Z[\theta]]$, but Theorem~\ref{thm:completions}
needs no hypothesis at all. Can you see how Hensel's factorisation lemma turns the Kummer--Dedekind factorisation into the
$p$-adic one when the $\bar g_i$ are distinct and $e_i=1$?
\end{think}

\section{Exercises}

\begin{exercise}[$\star$]\label{ex:quadQ3}
List the quadratic extensions of $\Q_3$ and of $\Q_5$, and say which are ramified.
\end{exercise}

\begin{exercise}[$\star$]\label{ex:twoplaces}
In $K=\Q(i)$, $5=(2+i)(2-i)$. Describe the two completions of $K$ above $5$: in each, which square root of $-1$ in $\Q_5$
does $i$ become? Compute $|2+i|$ in both.
\end{exercise}

\begin{exercise}[$\star\star$]\label{ex:sqrt2Q7}
Show that $\sqrt2\in\Q_7$, compute its first four $7$-adic digits by Newton's method, and describe the places of $\Q(\sqrt2)$
above $7$.
\end{exercise}

\begin{exercise}[$\star\star$]\label{ex:zetapEis}
Let $p$ be odd. Show that $\Q_p(\zeta_p)=\Q_p\bigl(\sqrt[p-1]{-p}\bigr)$.
\end{exercise}

\begin{exercise}[$\star\star$]\label{ex:seven}
\emph{The seven quadratic extensions of $\Q_2$.} Show that $\Q_2(\sqrt5)=\Q_2(\sqrt{-3})=\Q_2(\zeta_3)$ is unramified, and that
the six others are totally ramified, by exhibiting Eisenstein polynomials.
\end{exercise}

\begin{exercise}[$\star\star\star$]\label{ex:tame}
\emph{Tame extensions.} Let $L/\Q_p$ be totally ramified of degree $e$ with $p\nmid e$. Show that $L=\Q_p\bigl(\sqrt[e]{p\omega}
\bigr)$ for a root of unity $\omega\in\Z_p$ of order dividing $p-1$. Why does the hypothesis $p\nmid e$ matter?
\end{exercise}

\begin{exercise}[$\star\star\star$]\label{ex:QpQq}
Show that $\Q_p$ and $\Q_q$ are not isomorphic as fields when $p\ne q$. (A field isomorphism fixes $\Q$; look at $1+p$.)
\end{exercise}

\begin{beauty}
A number field is a single, rigid object; its completions are many, flexible, and simple. Each $K_P$ is a tower of two
transparent pieces, an unramified layer made of roots of unity and a totally ramified layer made of an Eisenstein root, and
Krasner's lemma says that the whole world of $p$-adic fields is discrete: wiggle the polynomial a little and the field does
not change. The art of modern number theory is to study $K$ through all its completions at once, real and $p$-adic, on an
equal footing. The product formula is the first sign that they are bound together, Hasse's principle (Chapter~9) the
second, and the ad\`eles of Chapter~12 make the idea into a single space.
\end{beauty}

\section{Hints}

Try each exercise for a while before reading its hint, and read the hint before the solution.

\begin{hint}{ex:quadQ3}
Use the square classes $1,u,p,up$.
\end{hint}
\begin{hint}{ex:twoplaces}
In the completion at $(2+i)$, the image of $2+i$ is small.
\end{hint}
\begin{hint}{ex:sqrt2Q7}
$3^2=9\equiv2\pmod 7$, and Newton's iteration is $x\mapsto\frac{x^2+2}{2x}$.
\end{hint}
\begin{hint}{ex:zetapEis}
Show $-p=(1-\zeta)^{p-1}w$ with $w\equiv1\pmod{1-\zeta}$, using Wilson's theorem, and take a $(p-1)$-st root of $w$ by Hensel.
\end{hint}
\begin{hint}{ex:seven}
Substitute $X\mapsto X+1$ where needed, and recall that $\frac{1+\sqrt5}2$ is a root of $X^2-X-1$.
\end{hint}
\begin{hint}{ex:tame}
Write $\pi_L^e=pu$, split off the Teichm\"uller representative of $u$, and take an $e$-th root of the rest by Hensel.
\end{hint}
\begin{hint}{ex:QpQq}
In $\Q_p$, $1+p$ is an $n$-th power for every $n$ prime to $p$. In $\Q_q$, what are the units that are $q^k$-th powers for every $k$?
\end{hint}

\section{Solutions}

\begin{solution}{ex:quadQ3}
\idea Quadratic extensions are square classes, and the class decides the ramification: a unit gives an unramified field,
$p$ times a unit an Eisenstein one.
\writeup For $p=3$, $u=-1$ is a non-square unit ($-1$ is not a square modulo $3$), so the extensions are $\Q_3(\sqrt{-1})$,
unramified ($X^2+1$ is irreducible modulo $3$, so the residue field is $\mathbb F_9$), and $\Q_3(\sqrt3)$, $\Q_3(\sqrt{-3})$, totally
ramified ($X^2\mp3$ is Eisenstein). For $p=5$, $u=2$: $\Q_5(\sqrt2)$ is unramified, and $\Q_5(\sqrt5)$, $\Q_5(\sqrt{10})$ are totally
ramified. Note that $\Q_5(\sqrt{-1})=\Q_5$, since $-1$ is a square modulo $5$.
\end{solution}

\begin{solution}{ex:twoplaces}
\idea The completion at $P$ is the embedding in which the elements of $P$ become $5$-adically small.
\writeup $X^2+1$ factors over $\Q_5$ as $(X-\iota)(X+\iota)$ with $\iota=2+5+2\cdot5^2+5^3+\dots\equiv2$ and
$-\iota=3+3\cdot5+2\cdot5^2+\dots\equiv3\pmod5$. So $K$ has two completions, both equal to $\Q_5$, given by $i\mapsto\iota$ and $i\mapsto
-\iota$. Under $i\mapsto-\iota$, $2+i\mapsto2-\iota\equiv0\pmod 5$, and in fact $|2-\iota|_5=5^{-1}$ (because $N(2+i)=5$ and
$|2+\iota|_5=1$). So this is the completion at $P=(2+i)$, where $|2+i|=\frac15$, while at $P'=(2-i)$, $i\mapsto\iota$ and
$|2+i|=|2+\iota|_5=1$. The product formula checks out: $|2+i|_\infty^2=5$, and $5\cdot\frac15\cdot1=1$.
\end{solution}

\begin{solution}{ex:sqrt2Q7}
\idea Hensel: a simple root modulo $7$ lifts uniquely, and Newton doubles the number of correct digits at each step.
\writeup $f=X^2-2$, $a=3$: $|f(3)|_7=|7|_7<1=|f'(3)|_7^2$, so $\sqrt2\in\Z_7$ with $\sqrt2\equiv3$. Newton: $x_1=\frac{9+2}{6}=\frac{11}6$;
modulo $7^2$, $6^{-1}\equiv41$, so $x_1\equiv11\cdot41=451\equiv451-9\cdot49=10=3+1\cdot7$. Next $x_2=\frac{x_1^2+2}{2x_1}$ modulo
$7^4$ gives $3+1\cdot7+2\cdot7^2+6\cdot7^3$ (check: $108^2-2=11662=7^3\cdot34$, so $108=3+7+2\cdot49$ is right modulo $7^3$). So
$\sqrt2=3+7+2\cdot7^2+6\cdot7^3+\dots$, and the other root is $-\sqrt2=4+5\cdot7+4\cdot7^2+0\cdot7^3+\dots$. In $K=\Q(\sqrt2)$,
$X^2-2$ splits into two distinct factors over $\Q_7$, so by Theorem~\ref{thm:completions} there are two primes above $7$, both
with $e=f=1$, and both completions are $\Q_7$, via $\sqrt2\mapsto\pm(3+7+2\cdot7^2+\dots)$. Concretely $7=(3+\sqrt2)(3-\sqrt2)$.
\end{solution}

\begin{solution}{ex:zetapEis}
\idea Both fields are totally ramified of degree $p-1$; it suffices to find a $(p-1)$-st root of $-p$ inside $\Q_p(\zeta)$,
and $1-\zeta$ is one up to a unit that is $\equiv1$, hence a $(p-1)$-st power.
\writeup Let $\lambda=1-\zeta$ and $L=\Q_p(\zeta)$. Since $\Phi_p(X+1)$ is Eisenstein at $p$ (Chapter~5), it is irreducible over
$\Q_p$ (Theorem~\ref{thm:eisensteinlocal}), so $[L:\Q_p]=p-1$ and $L$ is totally ramified with uniformiser $\lambda$. Now
$p=\prod_{a=1}^{p-1}(1-\zeta^a)=\lambda^{p-1}\prod_au_a$ with $u_a=1+\zeta+\dots+\zeta^{a-1}\equiv a\pmod\lambda$ (as $\zeta\equiv1$).
By Wilson's theorem $\prod_au_a\equiv(p-1)!\equiv-1$, so $-p=\lambda^{p-1}w$ with $w\equiv1\pmod\lambda$. For $g=X^{p-1}-w$:
$|g(1)|<1=|g'(1)|^2$ ($g'(1)=p-1$ is a unit), so Hensel in $\mathcal O_L$ gives $c$ with $c^{p-1}=w$. Then $\beta=\lambda c\in L$ has
$\beta^{p-1}=-p$. The polynomial $X^{p-1}+p$ is Eisenstein, so $\Q_p(\beta)$ has degree $p-1$, and $\Q_p(\beta)=L$.
\end{solution}

\begin{solution}{ex:seven}
\idea Unramified means the residue field grows; Eisenstein means a uniformiser of absolute value $2^{-1/2}$.
\writeup $\frac{1+\sqrt5}2$ is a root of $X^2-X-1$, irreducible modulo $2$, so $\Q_2(\sqrt5)$ has residue field $\mathbb F_4$ and is
unramified; by Theorem~\ref{thm:unram} it is $\Q_2(\zeta_3)=\Q_2(\sqrt{-3})$ (and indeed $(-3)\cdot5=-15\equiv1\pmod8$ is a square, so
$-3$ and $5$ lie in the same square class).
For the others: $X^2-2$, $X^2+2$, $X^2-10$, $X^2+10$ are Eisenstein; $\sqrt{-1}=1+\pi$ with $\pi$ a root of $(X+1)^2+1=X^2+2X+2$,
Eisenstein; $\sqrt{-5}=1+\pi$ with $\pi$ a root of $(X+1)^2+5=X^2+2X+6$, Eisenstein since $v_2(6)=1$. So these six are totally
ramified. In each case $\mathcal O_L=\Z_2[\pi]=\Z_2[\sqrt d]$ (Theorem~\ref{thm:efn} with $f=1$ and $\omega_1=1$), so the
discriminant is $\operatorname{disc}(X^2-d)=4d$: its $2$-adic valuation is $2$ for $d=-1,-5$ and $3$ for $d=\pm2,\pm10$, as in
Figure~\ref{fig:quadratic2}.
\end{solution}

\begin{solution}{ex:tame}
\idea Everything that is $\equiv1$ modulo the maximal ideal is an $e$-th power when $p\nmid e$, so only the residue of the unit
in $\pi_L^e=pu$ matters, and it is a root of unity.
\writeup Let $\pi=\pi_L$; $|\pi^e|=|p|$, so $\pi^e=pu$ with $u\in\mathcal O_L^\times$. Since $L/\Q_p$ is totally ramified, $k_L=\mathbb F_p$,
so $u\equiv a\pmod\pi$ with $a\in\{1,\dots,p-1\}$. Let $\omega\in\Z_p$ be the root of $X^{p-1}-1$ with $\omega\equiv a$ (Volume~1,
Corollary~12.9) and $u_1=u/\omega\equiv1\pmod\pi$. Hensel in $\mathcal O_L$ for $X^e-u_1$ at $1$ ($g'(1)=e$ is a unit since $p\nmid e$)
gives $v$ with $v^e=u_1$. Then $\beta=\pi/v$ satisfies $\beta^e=p\omega$. As $X^e-p\omega$ is Eisenstein, $\Q_p(\beta)$ has degree $e$,
so $L=\Q_p(\beta)$. If $p\mid e$, units $\equiv1$ need not be $e$-th powers: $\Q_2(\sqrt{-1})$ and $\Q_2(\sqrt2)$ are both totally
ramified of degree $2$, but only one of them has the form $\Q_2(\sqrt{2\omega})$ with $\omega\in\mu_1=\{1\}$.
\end{solution}

\begin{solution}{ex:QpQq}
\idea An isomorphism is the identity on $\Q$, so it preserves which rational numbers have $n$-th roots. In $\Q_p$, the number
$1+p$ has roots of all orders prime to $p$; in $\Q_q$, a rational number with $q^k$-th roots for all $k$ is forced to be a root
of unity.
\writeup Any field isomorphism $\phi:\Q_p\to\Q_q$ fixes $1$, hence $\Q$. One of $p,q$ is odd; by symmetry say $q$ is odd. In $\Q_p$,
$x=1+p$ is an $n$-th power for every $n$ prime to $p$: Hensel for $X^n-x$ at $1$, as $|1-x|_p<1=|n|_p^2$. In particular, for every
$k$, $1+p=y^{q^k}$ in $\Q_p$, hence in $\Q_q$ (apply $\phi$). In $\Q_q$, $|y|_q=1$, and $z=y^{q-1}\equiv1\pmod q$ by Fermat in
$\mathbb F_q$; by the binomial theorem $(1+qt)^{q^k}\equiv1\pmod{q^{k+1}}$ for $t\in\Z_q$. So $(1+p)^{q-1}=z^{q^k}\equiv1\pmod
{q^{k+1}}$ for all $k$, that is, $(1+p)^{q-1}=1$. But $(1+p)^{q-1}>1$ is a rational number. Contradiction.
\end{solution}

\section*{Sources and further reading}

\begin{itemize}
\item K.~Hensel, ``\"Uber eine neue Begr\"undung der Theorie der algebraischen Zahlen'', \emph{Jahresber.\ Deutsch.\ Math.-Verein.}\
6 (1897), 83--88; and \emph{Theorie der algebraischen Zahlen}, Teubner, 1908. The birth of $p$-adic numbers.
\item F.~Q.~Gouv\^ea, \emph{$p$-adic Numbers: An Introduction}, 2nd ed., Springer, 1997. A gentle book; Chapter~5 on finite
extensions and $\C_p$, and historical notes including Hensel's mistaken proof.
\item J.~Neukirch, \emph{Algebraic Number Theory}, Springer, 1999, Chapter~II, \S\S3--8. Hensel's lemma, extensions of
valuations, unramified and tamely ramified extensions, and Krasner's lemma.
\item J.-P.~Serre, \emph{Local Fields}, Springer GTM~67, 1979, Chapters~I--III; and ``Une `formule de masse' pour les extensions
totalement ramifi\'ees de degr\'e donn\'e d'un corps local'', \emph{C.\ R.\ Acad.\ Sci.\ Paris} 286 (1978), 1031--1036.
\item J.~W.~S.~Cassels, \emph{Local Fields}, Cambridge University Press, 1986. Many worked examples, and Ostrowski's theorems.
\item M.~Krasner, ``Th\'eorie non ab\'elienne des corps de classes pour les extensions finies et s\'eparables des corps valu\'es
complets'', \emph{C.\ R.\ Acad.\ Sci.\ Paris} 222 (1946), 626--628. The lemma.
\item H.~Hasse, ``\"Uber die Darstellbarkeit von Zahlen durch quadratische Formen im K\"orper der rationalen Zahlen'', \emph{J.\
reine angew.\ Math.}\ 152 (1923), 129--148. The local--global principle, for Chapter~9.
\item \emph{History.} Hensel's flawed 1905 proof and Hasse's discovery of Hensel's book in a second-hand bookshop are told in
Gouv\^ea's book above (historical notes). K\"ursch\'ak (1913) and Ostrowski (1916): P.~Roquette, ``History of valuation
theory, Part~I'', in \emph{Valuation Theory and Its Applications}, vol.~I, Fields Inst.\ Commun.\ 32 (2002), 291--355.
\end{itemize}
